% !TEX root = chapter-standalone.tex

\devel{
\section[Questions and answers]{Questions and answers}
\label{sec:modeling-cats-qa}

\begin{ctdefinition}[GC]\label{def:GC-definition}
        An object of $\mathbf{GSet}$ is a tuple
        \begin{equation}
            \tup{Q, A, R},
        \end{equation}
        where $Q$ and $A$ are sets, and $R : Q \toinRel A$ is a relation.

        A morphism
        $
            \mathbf{r} \colon \tupp{Q_1, A_1, R_1}
            \mtoin{\mathbf{GSet}}
            \tupp{Q_2, A_2, R_2}
        $
        is a pair of maps
        \begin{align}
            \mathbf{r} = & \,\tupp{r_\flat, r^\sharp}, \\
                         & \ r_\flat: Q_1 \mfromin\Set Q_2, \\
                         & \ r^\sharp: A_1 \mtoin\Set A_2,
        \end{align}
        that satisfy the property
        \begin{equation}\label{eq:gset-mor-cond}
            \forall q_2: Q_2 \quad \forall a_1: A_1 \quad
            \inrel{r_\flat(q_2)}{R_1}{a_1}
            \ \Imp\
            \inrel{q_2}{R_2}{r^\sharp(a_1)}.
        \end{equation}
        Morphism composition is defined component-wise
        \begin{align}
            ( \mathbf{r} \mthen \mathbf{s} )_\flat  & = s_\flat \mthen r_\flat, \\
            ( \mathbf{r} \mthen \mathbf{s} )^\sharp & = r^\sharp \mthen s^\sharp.
        \end{align}

        The identity at $\tup{Q, A, R}$ is
        $
            \catidat{\tup{Q, A, R}} =
            \tup{
                \catidat{Q}, \catidat{A}
            }.
        $
    \end{ctdefinition}


\begin{gradedexercise}[\exname{HwkGSet}]
    \label{ex:HwkGSet}

  Check that the definition of $\mathbf{GSet}$ does indeed define a category.

    \end{gradedexercise}

\solutionof{HwkGSet}

}
